% Generated from locale/tr/content/many-valued-logic/sequent-calculus/propositional-rules.tex; do not hand-edit.


\olfileid[tr]{mvl}{seq}{prl}

\olsection{Seçilmiş Mantıklar İçin Önerme Kuralları}

Bir $n$-yanlı sekant hesabında bir bağlaca ilişkin çıkarım kuralları
yalnızca o bağlacın karakteristik doğruluk fonksiyonuna bağlıdır.
Dolayısıyla bir bağlaç farklı mantıklarda aynı doğruluk fonksiyonuyla
tanımlanıyorsa, bağlacın $n$-yanlı sekant kuralları bu mantıklarda aynıdır.

\subsection{$\lnot$ İçin Kurallar}

Aşağıdaki $\lnot$ kuralları \L ukasiewicz ve Kleene mantıkları ile bunların
çeşitlemelerinde geçerlidir.

\begin{defish}
  \begin{center}
\AxiomC{$ \Gamma \nSequent \Pi \nSequent \Delta, !A $}
\RightLabel{\iR{\lnot}{\False}}
\UnaryInfC{$ \lnot !A, \Gamma \nSequent \Pi \nSequent \Delta$}
\DisplayProof
\\[2ex]
\AxiomC{$ \Gamma \nSequent !A, \Pi \nSequent \Delta $}
\RightLabel{\iR{\lnot}{\Undef}}
\UnaryInfC{$\Gamma \nSequent \lnot !A, \Pi \nSequent \Delta$}
\DisplayProof
\\[2ex]
\AxiomC{$!A, \Gamma \nSequent \Pi \nSequent \Delta$}
\RightLabel{\iR{\lnot}{\True}}
\UnaryInfC{$\Gamma \nSequent \Pi \nSequent \Delta,  \lnot !A$}
\DisplayProof
  \end{center}
\end{defish}

Aşağıdaki $\lnot$ kuralları G\"odel mantığında geçerlidir.

\begin{defish}
\AxiomC{$ \Gamma \nSequent !A, \Pi \nSequent \Delta, !A $}
\RightLabel{\iR{\lnot}{\False}[\LogGod]}
\UnaryInfC{$ \lnot !A, \Gamma \nSequent \Pi \nSequent \Delta$}
\DisplayProof
\hfill
\AxiomC{$!A, \Gamma \nSequent \Pi \nSequent \Delta$}
\RightLabel{\iR{\lnot}{\True}[\LogGod]}
\UnaryInfC{$\Gamma \nSequent \Pi \nSequent \Delta,  \lnot !A$}
\DisplayProof
\hfill
\end{defish}

(G\"odel mantığında $\lnot !A$ hiçbir zaman $\Undef$ değerini alamaz;
dolayısıyla orta konum için bir kural yoktur.)

\subsection{$\land$ İçin Kurallar}

Bunlar \L ukasiewicz, güçlü Kleene ve G\"odel mantıklarında $\land$ için
geçerli kurallardır.

\begin{defish}
\begin{center}
  \AxiomC{$!A, !B, \Gamma \nSequent \Pi \nSequent \Delta$}
  \RightLabel{$\iR{\land}{\False}$}
  \UnaryInfC{$!A \land !B, \Gamma \nSequent \Pi \nSequent \Delta$}
  \DisplayProof
\\[2ex]
  \AxiomC{$\Gamma \nSequent !A, \Pi \nSequent !A, \Delta$}
  \AxiomC{$\Gamma \nSequent !B, \Pi \nSequent !B, \Delta$}
  \AxiomC{$\Gamma \nSequent !A, !B, \Pi \nSequent \Delta$}
  \RightLabel{$\iR\land\Undef$}
  \TrinaryInfC{$\Gamma \nSequent !A \land !B, \Pi \nSequent \Delta$}
  \DisplayProof
  \\[2ex]  
  \AxiomC{$\Gamma \nSequent \Pi \nSequent \Delta, !A$}
  \AxiomC{$\Gamma \nSequent \Pi \nSequent \Delta, !B$}
  \RightLabel{$\iR\land\True$}
  \BinaryInfC{$\Gamma \nSequent \Pi \nSequent \Delta, !A \land !B$}
  \DisplayProof
\end{center}
\end{defish}

\subsection{$\lor$ İçin Kurallar}

Bunlar \L ukasiewicz, güçlü Kleene ve G\"odel mantıklarında $\lor$ için
geçerli kurallardır.

\begin{defish}
\begin{center}
 \AxiomC{$!A, \Gamma \nSequent \Pi \nSequent \Delta$}
  \AxiomC{$!B, \Gamma \nSequent \Pi \nSequent \Delta$}
  \RightLabel{$\iR\lor\False$}
  \BinaryInfC{$!A \lor !B, \Gamma \nSequent \Pi \nSequent \Delta$}
  \DisplayProof
\\[2ex]
  \AxiomC{$!A, \Gamma \nSequent !A, \Pi \nSequent \Delta$}
  \AxiomC{$!B, \Gamma \nSequent !B, \Pi \nSequent \Delta$}
  \AxiomC{$\Gamma \nSequent !A, !B, \Pi \nSequent \Delta$}
  \RightLabel{$\iR\lor\Undef$}
  \TrinaryInfC{$\Gamma \nSequent !A \lor !B, \Pi \nSequent \Delta$}
  \DisplayProof
  \\[2ex]  
  \AxiomC{$\Gamma \nSequent \Pi \nSequent \Delta, !A, !B$}
  \RightLabel{$\iR\lor\True$}
  \UnaryInfC{$\Gamma \nSequent \Pi \nSequent \Delta, !A \lor !B$}
  \DisplayProof
\end{center}
\end{defish}

\subsection{$\lif$ İçin Kurallar}

Bunlar \L ukasiewicz mantığında $\lif$ için geçerli kurallardır.

\begin{defish}
\begin{center}
  \AxiomC{$\Gamma \nSequent \Pi \nSequent \Delta, !A$}
  \AxiomC{$!B, \Gamma \nSequent \Pi \nSequent \Delta$}
  \RightLabel{$\iR\lif\False[\LogLuk[3]]$}
  \BinaryInfC{$!A \lif !B, \Gamma \nSequent \Pi \nSequent \Delta$}
  \DisplayProof
  \\[2ex]
  \AxiomC{$\Gamma \nSequent !A, !B, \Pi \nSequent \Delta$}
  \AxiomC{$!B, \Gamma \nSequent \Pi \nSequent \Delta, !A$}
  \RightLabel{$\iR\lif\Undef[\LogLuk[3]]$}
  \BinaryInfC{$\Gamma \nSequent !A \lif !B, \Pi \nSequent \Delta$}
  \DisplayProof
\\[2ex]
  \AxiomC{$!A, \Gamma \nSequent !B, \Pi \nSequent \Delta, !B$}
  \AxiomC{$!A, \Gamma \nSequent !A, \Pi \nSequent \Delta, !B$}
  \RightLabel{$\iR{\lif}{\True}[\LogLuk[3]]$}
  \BinaryInfC{$\Gamma \nSequent \Pi \nSequent \Delta, !A \lif !B$}
  \DisplayProof
\end{center}
\end{defish}

Bunlar güçlü Kleene mantığında $\lif$ için geçerli kurallardır.

\begin{defish}
\begin{center}
  \AxiomC{$\Gamma \nSequent \Pi \nSequent \Delta, !A$}
  \AxiomC{$!B, \Gamma \nSequent \Pi \nSequent \Delta$}
  \RightLabel{$\iR\lif\False[\LogKs]$}
  \BinaryInfC{$!A \lif !B, \Gamma \nSequent \Pi \nSequent \Delta$}
  \DisplayProof
  \\[2ex]
  \AxiomC{$!B, \Gamma \nSequent !B, \Pi \nSequent \Delta$}
  \AxiomC{$\Gamma \nSequent !A, !B, \Pi \nSequent \Delta$}
  \AxiomC{$\Gamma \nSequent !A, \Pi \nSequent \Delta, !A$}
  \RightLabel{$\iR\lif\Undef[\LogKs]$}
  \TrinaryInfC{$\Gamma \nSequent !A \lif !B, \Pi \nSequent \Delta$}
  \DisplayProof
\\[2ex]
  \AxiomC{$!A, \Gamma \nSequent \Pi \nSequent \Delta, !B$}
  \RightLabel{$\iR{\lif}{\True}[\LogKs]$}
  \UnaryInfC{$\Gamma \nSequent \Pi \nSequent \Delta, !A \lif !B$}
  \DisplayProof
\end{center}
\end{defish}

Bunlar G\"odel mantığında $\lif$ için geçerli kurallardır.

\begin{defish}
\begin{center}
  \AxiomC{$\Gamma \nSequent !A, \Pi \nSequent \Delta, !A$}
  \AxiomC{$!B, \Gamma \nSequent \Pi \nSequent \Delta$}
  \RightLabel{$\iR\lif\False[\LogGod[3]]$}
  \BinaryInfC{$!A \lif !B, \Gamma \nSequent \Pi \nSequent \Delta$}
  \DisplayProof
  \\[2ex]
  \AxiomC{$\Gamma \nSequent !B, \Pi \nSequent \Delta$}
  \AxiomC{$\Gamma \nSequent \Pi \nSequent \Delta, !A$}
  \RightLabel{$\iR\lif\Undef[\LogGod[3]]$}
  \BinaryInfC{$\Gamma \nSequent !A \lif !B, \Pi \nSequent \Delta$}
  \DisplayProof
\\[2ex]
  \AxiomC{$!A, \Gamma \nSequent !B, \Pi \nSequent \Delta, !B$}
  \AxiomC{$!A, \Gamma \nSequent !A, \Pi \nSequent \Delta, !B$}
  \RightLabel{$\iR{\lif}{\True}[\LogGod[3]]$}
  \BinaryInfC{$\Gamma \nSequent \Pi \nSequent \Delta, !A \lif !B$}
  \DisplayProof
\end{center}
\end{defish}


\begin{sidewaysfigure}
  \begin{center}  
  \AxiomC{$A \nSequent A \nSequent A$}
  \RightLabel{\iR \Weakening \True}
  \UnaryInfC{$A \nSequent A \nSequent B, A$}
  \RightLabel{\iR \Weakening \Undef}
  \UnaryInfC{$A \nSequent B, A \nSequent B, A$}
  \RightLabel{\iR \Weakening \Undef}
  \UnaryInfC{$A \nSequent A, B, A \nSequent B, A$}
  \AxiomC{$A \nSequent A \nSequent A$}
  \RightLabel{\iR \Weakening \True}
  \UnaryInfC{$A \nSequent A \nSequent A, A$}
  \RightLabel{\iR \Weakening \True}
  \UnaryInfC{$A \nSequent A \nSequent B, A, A$}
  \RightLabel{\iR \Weakening \False}
  \UnaryInfC{$B, A \nSequent A \nSequent B, A, A$}
  \RightLabel{$\iR\lif\Undef$}
  \BinaryInfC{$A \nSequent A \lif B, A \nSequent B, A$}
  \AxiomC{$B \nSequent B \nSequent B$}
  \RightLabel{\iR \Weakening \Undef}
  \UnaryInfC{$B \nSequent A, B \nSequent B$}
  \RightLabel{\iR \Exchange \Undef}
  \UnaryInfC{$B \nSequent B, A \nSequent B$}
  \RightLabel{\iR \Weakening \Undef}
  \UnaryInfC{$B \nSequent A, B, A \nSequent B$}
  \RightLabel{\iR \Weakening \False}
  \UnaryInfC{$A, B \nSequent A, B, A \nSequent B$}
  \RightLabel{\iR \Exchange \False}
  \UnaryInfC{$B, A \nSequent A, B, A \nSequent B$}
  \AxiomC{$A \nSequent A \nSequent A$}
  \RightLabel{\iR \Weakening \True}
  \UnaryInfC{$A \nSequent A \nSequent B, A$}
  \RightLabel{\iR \Weakening \False}
  \UnaryInfC{$B, A \nSequent A \nSequent B, A$}
  \RightLabel{\iR \Weakening \False}
  \UnaryInfC{$B, B, A \nSequent A \nSequent B, A$}
  \RightLabel{$\iR\lif\Undef$}
  \BinaryInfC{$B, A \nSequent A \lif B, A \nSequent B$}
  \RightLabel{$\iR\lif\False$}
  \BinaryInfC{$A \lif B, A \nSequent A \lif B, A \nSequent B$}
  \DisplayProof
  \end{center}
  \caption{$\LogLuk[3]$ İçinde Örnek Bir !!{derivation}}
\end{sidewaysfigure}

